\documentclass[10pt,a4paper]{article}
\usepackage[margin=22mm]{geometry}
\usepackage{amsmath,amssymb,amsthm,mathtools}
\usepackage[T1]{fontenc}
\usepackage{lmodern}
\usepackage[hidelinks]{hyperref}
\setlength{\parindent}{0pt}
\setlength{\parskip}{5pt}
\newcommand{\E}{\mathbb E}
\newcommand{\Var}{\operatorname{Var}}
\newcommand{\Pbb}{\mathbb P}
\newcommand{\R}{\mathbb R}
\newtheorem{theorem}{Theorem}
\title{A one-dimensional obstruction to the printed\protect\\random-polytope central limit assertion}
\author{lizaixi01 \quad Problem 00000001453}
\date{24 October 2026}
\begin{document}
\maketitle
\vspace{-7mm}

\textbf{Scope.} The source defines $K_n$ as the convex hull of uniform points in a
convex body $K\subset\R^d$, and asserts that the normalized volume deficit obeys a
central limit theorem. Both printed versions omit a restriction $d\ge2$.
We take the usual positive-dimensional convex-body convention and the usual
centering by the actual expectation and scaling by the actual standard deviation.
The counterexample is the full-dimensional body $K=[0,1]\subset\R$, not a
lower-dimensional subset of a higher-dimensional space. Only the CLT conjunct is
refuted. No value is assigned to the source's unspecified sphere constants, and
no contradiction to a corrected $d\ge2$ theorem is claimed.

\begin{theorem}
Let $n\ge2$, let $X_1,\ldots,X_n$ be independent uniform samples in $[0,1]$,
and let $Y_n=\operatorname{vol}(K)-\operatorname{vol}(K_n)$. Then
\begin{align*}
\E Y_n&=\frac2{n+1},&
\E Y_n^2&=\frac6{(n+1)(n+2)},&
\Var(Y_n)&=\frac{2(n-1)}{(n+1)^2(n+2)}>0.
\end{align*}
The actual standardized laws of
$Z_n=(Y_n-\E Y_n)/\sqrt{\Var(Y_n)}$ do not converge weakly to $N(0,1)$;
their CDFs also fail the ordinary Gaussian CLT criterion.
\end{theorem}

\textbf{The true sampling model and hull.}
The measure $U=\lambda|_{[0,1]}$ is already normalized, since
$\lambda([0,1])=1$. The sample measure is the genuine finite product $U^{\otimes n}$,
with coordinate projections as the random variables. Thus the coordinates are
independent and have the normalized restriction law. The interval is compact,
convex and has nonempty interior, and $\dim_\R\R=1$.
Put $A=\min_iX_i$, $M=\max_iX_i$ and $B=1-M$. For every finite sample the true
convex hull is $[A,M]$: all points lie in this interval, and both endpoints
are sample points. Its Lebesgue volume is $M-A$. Consequently
$Y_n=A+B$, with $A,B\in[0,1]$ almost surely.

\textbf{Moments derived from the product measure.}
For $u,v\in[0,1]$, the coordinate rectangle gives exactly
\[
 \Pbb(A\ge u,B\ge v)=\bigl(\max(1-u-v,0)\bigr)^n,
 \qquad
 \Pbb(A\ge u)=\Pbb(B\ge u)=(1-u)^n.
\]
The joint formula includes the genuinely empty rectangle when $u+v>1$;
it is proved from the measure of that rectangle, not postulated as an
order-statistic distribution.

For a measurable $X\in[0,1]$ and bounded measurable $W$ on any finite measure
space, Fubini applied to
$\mathbf1_{t\le X(\omega)}W(\omega)t^k$ proves
\[
 \int W\,\frac{X^{k+1}}{k+1}
 =\int_0^1 t^k\left(\int_{\{X\ge t\}}W\right)dt.
\]
The Lean proof first verifies the pointwise cut integral
$\int_0^1\mathbf1_{t\le a}t^kdt=a^{k+1}/(k+1)$ and integrability on the
actual product measure. Taking $W=1$, $k=0,1$ and using the marginal tails yields
\[
 \E A=\E B=\frac1{n+1},\qquad
 \E A^2=\E B^2=\frac2{(n+1)(n+2)}.
\]
Applying the same identity with $X=A,W=B$, then to $B$ under the restricted
measure on $\{A\ge u\}$, yields the true cross moment
\[
\begin{aligned}
 \E(AB)
 &=\int_0^1\int_0^1\bigl(\max(1-u-v,0)\bigr)^n\,dv\,du\\
 &=\int_0^1\frac{(1-u)^{n+1}}{n+1}\,du
 =\frac1{(n+1)(n+2)}.
\end{aligned}
\]
Thus expanding $(A+B)^2$ proves the two displayed moments of $Y_n$.
The native variance identity $\Var(Y_n)=\E(Y_n^2)-(\E Y_n)^2$ gives the displayed
positive variance. All these quantities are integrals of the actual hull deficit.

\textbf{An obstruction at a Gaussian continuity point.}
Nonnegativity of $Y_n$ and the actual moment formulas imply
\[
 \frac{\E Y_n}{\sqrt{\Var(Y_n)}}
 =\sqrt{\frac{2(n+2)}{n-1}}\le\sqrt8<3,
 \qquad Z_n\ge-3\quad\hbox{almost surely}.
\]
In particular, every actual pushforward law $\nu_n$ satisfies
$\nu_n(({-\infty},-4])=0$. The genuine standard Gaussian measure $\gamma$ has
$\gamma(({-\infty},-4])>0$, because its density is strictly positive and
$[-5,-4]$ has positive Lebesgue measure. Its singleton boundary mass at $-4$ is zero.
The native Portmanteau theorem would therefore force
$\nu_n(({-\infty},-4])\to\gamma(({-\infty},-4])$ under weak convergence, an
impossibility. The same point directly contradicts convergence of the real-valued
CDFs. This proves both failures in the theorem.

\textbf{Formalization and reproduction.}
The Lean family uses $n=k+2$ with $k\in\mathbb N$, a cofinal indexing of all sample
sizes at least two. In \texttt{Main.lean}, \texttt{polytope} is the native convex
hull; \texttt{deficit} uses its native Lebesgue measure. The weighted-tail theorem
proves the Fubini identity; the endpoint and deficit moment theorems compute the
actual integrals and native variance. \texttt{standardizedLaw} is the native
probability-measure pushforward, and the final weak-convergence and CDF theorems
prove the convergence obstructions. There is no assumed density identification
or external geometric identification. Literature is used only for the customary
model and normalization, not as a proof of this counterexample.

The project pins Lean 4.33.0 and Mathlib commit
\begin{center}\texttt{db584cd6d46c92f209a44c0f1c829460d327499d}.\end{center}
From \texttt{lean/}, run \texttt{lean --version}, then \texttt{./verify.sh}.
The script runs the following build, exhaustive ordinary axiom check and fresh check:
\begin{verbatim}
lake build +Main
lake env lean -DwarningAsError=true Check.lean
lake env leanchecker --fresh Main
\end{verbatim}
Only \texttt{propext}, \texttt{Classical.choice} and \texttt{Quot.sound} are allowed.
No admitted proof, native evaluator proof or additional axiom is used.

\begin{thebibliography}{9}
\bibitem{vu} V. Vu, \emph{Central limit theorems for random polytopes in a smooth
convex set}, arXiv:math/0503559v1 (2005), pp.~1--2 (model and standardized CLT),
p.~4 (variance order). \url{https://arxiv.org/abs/math/0503559}.
\bibitem{chen} Y.-C. Chen, \emph{Order statistics from a continuous univariate
distribution}, UW STAT 512, Lecture 9 (2022 Autumn), uniform-distribution case.
\url{https://faculty.washington.edu/yenchic/20A_stat512/Lec9_Order.html}.
\end{thebibliography}
\end{document}
